% Independently written draft from Illusie I, 1.1.6, printed pp.4-6.
% Local labels only. Source snippet for bounded root-chapter composition.

\section{Relative simplicial homotopies and base change}
\label{section-illusie-I-relative-homotopy}

\noindent
The following construction makes explicit the relative homotopy formalism
of Illusie, \textit{Complexe cotangent et d\'eformations I}, I.1.1.6.
The underlying absolute notion is that of
Remark \ref{remark-homotopy-better}; strongly cartesian morphisms and fibred
categories are as in
\href{https://stacks.math.columbia.edu/tag/02XK}{Tag 02XK} and
\href{https://stacks.math.columbia.edu/tag/02XM}{Tag 02XM}.

Let $p:\mathcal E\to\mathcal B$ be a functor. Write
$\mathrm{Simp}(\mathcal E)$ for the category of functors
$\Delta^{\mathrm{op}}\to\mathcal E$. If $X,Y$ are simplicial objects,
a simplicial homotopy is a family
$$
h_{n,\alpha}:X_n\longrightarrow Y_n,
\qquad \alpha:[n]\longrightarrow[1],
$$
whose values at the two constant maps are its endpoints and such that
$$
Y(\theta)h_{n,\alpha}
=h_{m,\alpha\theta}X(\theta)
\quad\text{for every }\theta:[m]\longrightarrow[n].
$$
This is the component formulation of simplicial homotopy; no coproducts
in $\mathcal E$ are required. Suppose $pX=S$, $pY=T$, and that both endpoints
lie over $f:S\to T$. Call the homotopy \emph{relative to $f$} if
$p(h_{n,\alpha})=f_n$ for every $n,\alpha$.
On each set of arrows over a fixed $f$, take the equivalence relation
generated by these homotopies in either direction.

\begin{lemma}
\label{lemma-illusie-I-relative-homotopy-quotient}
The relative homotopy relations are compatible with composition. They define
a category $\mathrm{hSimp}(\mathcal E/\mathcal B)$ over
$\mathrm{Simp}(\mathcal B)$, with the same objects as
$\mathrm{Simp}(\mathcal E)$. Its fibre over $S$ has as morphisms the
$\mathrm{id}_S$-relative homotopy classes of arrows over $S$.
Every functor $F:\mathcal E\to\mathcal E'$ over $\mathcal B$ induces a
functor between these quotient categories over $\mathrm{Simp}(\mathcal B)$.
\end{lemma}

\begin{proof}
Let $h$ be relative to $f:S\to T$. Precomposition with an arrow $v$
over $g:S'\to S$ replaces its components by $h_{n,\alpha}v_n$;
their images under $p$ are $f_ng_n$. Postcomposition with an arrow $w$
over $k:T\to T'$ gives $w_nh_{n,\alpha}$ over $k_nf_n$.
The displayed naturality relation is preserved in each case. Thus
precomposition and postcomposition preserve elementary relative homotopies,
and hence their generated equivalence relations. The quotient has
well-defined associative composition and identities inherited from the
original category. Applying $F$ to every component preserves naturality,
endpoints and the base map, proving the last assertion.
\end{proof}

\begin{lemma}
\label{lemma-illusie-I-relative-homotopy-base-change}
Assume $p:\mathcal E\to\mathcal B$ is a fibred category. Then:
\begin{enumerate}
\item For every $f:S\to T$ in $\mathrm{Simp}(\mathcal B)$ and every
$Y$ over $T$, there is a lift $c:f^*Y\to Y$ over $f$ whose components are
strongly cartesian. For every $X$ over $S$ and pair of arrows $a,b:X\to Y$
over $f$, composition with $c$ bijects homotopies over $\mathrm{id}_S$
between the lifted arrows $\widetilde a,\widetilde b:X\to f^*Y$ with
homotopies relative to $f$ between $a,b$.
\item Every $g:S'\to S$ induces a pullback functor
$$
g^*:\mathrm{hSimp}(\mathcal E)_S
\longrightarrow\mathrm{hSimp}(\mathcal E)_{S'},
$$
where these are the fibres of the preceding quotient category. The
canonical comparisons between successive pullbacks are well-defined on
these fibres and satisfy the usual identity and associativity constraints.
\item The functor
$\mathrm{hSimp}(\mathcal E/\mathcal B)\to\mathrm{Simp}(\mathcal B)$
is fibred. The class of a componentwise strongly cartesian lift is strongly
cartesian in this quotient.
\end{enumerate}
\end{lemma}

\begin{proof}
Choose a strongly cartesian arrow $c_n:(f^*Y)_n\to Y_n$ over $f_n$
in every degree. For $\theta:[m]\to[n]$, define $(f^*Y)(\theta)$
as the unique arrow over $S(\theta)$ satisfying
$$
c_m(f^*Y)(\theta)=Y(\theta)c_n.
$$
It exists since $T(\theta)f_n=f_mS(\theta)$. Uniqueness proves the
identity and composition laws, so $f^*Y$ is a simplicial object and $c$
is a simplicial arrow. The same argument lifts every arrow to $Y$ over
$f\circ u$ uniquely through $c$ over $u$, proving that $c$ is strongly
cartesian in the simplicial category.

Given an $f$-relative homotopy $h$ between $a$ and $b$, lift each component
uniquely through $c_n$ over $\mathrm{id}_{S_n}$:
$$
c_n\widetilde h_{n,\alpha}=h_{n,\alpha}.
$$
To check naturality, compose
$(f^*Y)(\theta)\widetilde h_{n,\alpha}$ and
$\widetilde h_{m,\alpha\theta}X(\theta)$ with $c_m$.
Their composites agree by the naturality of $h$; they lie over the same
map $S(\theta)$, so strong cartesianness of $c_m$ gives equality.
Uniqueness also identifies the endpoints with $\widetilde a,\widetilde b$.
Conversely, composition with $c$ gives a relative homotopy. These two
constructions are inverse, proving (1).

For (2), choose lifts $c_X:g^*X\to X$ and $c_Y:g^*Y\to Y$.
An $S$-relative homotopy $h$ between arrows $X\to Y$ gives an
$S'$-relative homotopy by the unique componentwise solutions of
$$
(c_Y)_n(g^*h)_{n,\alpha}=h_{n,\alpha}(c_X)_n.
$$
The naturality and endpoint checks are the same as in (1). Pullback
therefore preserves elementary homotopies and their equivalence closure,
and hence induces the displayed functor. Different choices of lifts have
unique compatible isomorphisms. Those isomorphisms identify the lifted
homotopies by uniqueness as well. A composite of strongly cartesian
arrows is strongly cartesian; comparing a successive lift with a lift
of the composite base map thus gives the canonical comparison
isomorphism. Identity and associativity constraints follow because both
routes are isomorphisms with the same composite to the original object
over the same base map. They remain equal after passing to the quotient.

Finally, take an arrow $Z\to Y$ over $f\circ u$, where $pZ=R$ and
$u:R\to S$. The strong cartesian property of $c$ gives its unique lift
$Z\to f^*Y$ over $u$. The componentwise argument of (1), now with
$u_n$ in place of $\mathrm{id}_{S_n}$, gives a bijection on the
corresponding sets of elementary relative homotopies. Every finite chain
of such homotopies also lifts uniquely, and composition preserves every
such chain. Consequently the bijection on arrows induces a bijection on
relative homotopy classes over each $u$. This is exactly the strongly
cartesian universal property in the quotient, proving (3).
\end{proof}

\begin{remark}
\label{remark-illusie-I-relative-homotopy-base-functor}
A functor over $\mathcal B$ always induces the quotient functor of the
first lemma. If both categories over $\mathcal B$ are fibred and the
functor preserves strongly cartesian arrows, the images
of the chosen lifts provide the canonical comparison with base change,
and the induced quotient functor preserves strongly cartesian arrows.
Indeed, every strongly cartesian arrow in the quotient is a chosen
componentwise strongly cartesian lift composed with a vertical isomorphism.
The quotient functor preserves this isomorphism as well as the chosen lift.
The latter claim is not asserted for an arbitrary functor over the base.
For a cofibred category and $g:S\to S'$, choose componentwise strongly
cocartesian lifts $d_X:X\to g_!X$. A relative homotopy $h:X\to Y$
over $S$ lifts by the uniquely determined equations
$$
(g_!h)_{n,\alpha}(d_X)_n=(d_Y)_nh_{n,\alpha}.
$$
Uniqueness gives naturality and the endpoints. The same equations with an
arbitrary subsequent base map prove the quotient cocartesian universal
property, just as in the preceding proof. Thus the quotient is cofibred
and $g_!$ points from the fibre over $S$ to the fibre over $S'$.

Cosimplicial variants follow without new lifting assumptions. A
cosimplicial object of $\mathcal E$ is a simplicial object of
$\mathcal E^{\mathrm{op}}$ after taking opposites, and this correspondence
reverses morphisms. Use the dual homotopy notion of
Lemma \ref{lemma-compare-homotopies}. Passing to opposites changes
strongly cartesian arrows into strongly cocartesian arrows and conversely.
Apply the corresponding simplicial result over $\mathcal B^{\mathrm{op}}$
and then take opposites again. This produces the relative cosimplicial
quotient, its base-change functors, and their coherence, with fibred
and cofibred hypotheses in their original directions.
\end{remark}

\begin{definition}[Signed total of an $n$-complex]
\label{definition-illusie-I-signed-total}
Let $\mathcal A$ be additive. An $n$-complex is a $\mathbb Z^n$-graded
object $L$ with maps $d_i:L^{\mathbf p}\to L^{\mathbf p+\mathbf e_i}$
satisfying $d_i^2=0$ and $d_id_j=d_jd_i$. Suppose that, for every $r$,
only finitely many $L^{\mathbf p}$ with $|\mathbf p|=r$ are nonzero (or,
more generally, that the indicated finite or existing coproduct exists).
Its signed total is
\[
 \operatorname{Tot}(L)^r=\bigoplus_{|\mathbf p|=r}L^{\mathbf p},
 \qquad d|_{L^{\mathbf p}}=\sum_{j=1}^n
 (-1)^{p_1+\cdots+p_{j-1}}d_j.
\]
\end{definition}

\begin{lemma}[The total differential]
\label{lemma-illusie-I-signed-total-differential}
The signed maps in Definition
\ref{definition-illusie-I-signed-total} square to zero. Totalization is
functorial in maps of $n$-complexes and is unchanged by grouping adjacent
coordinates and totalizing again.
\end{lemma}

\begin{proof}
On $L^{\mathbf p}$, the two terms involving $d_jd_k$ with $j<k$ have
opposite signs: after the $j$th differential the prefix exponent for the
$k$th term has gained one. They therefore cancel because
$d_jd_k=d_kd_j$. The terms
$d_j^2$ vanish. The same calculation proves functoriality. For a grouping,
the exponent contributed by the outer differential is the sum of the
degrees preceding the grouped block, while the inner exponent is the sum
inside it; their sum is exactly the displayed prefix exponent. The direct
sum indexing is the same, so the two totalizations agree canonically.
\end{proof}

\begin{lemma}[Multisimplicial chains]
\label{lemma-illusie-I-multisimplicial-chains}
Let $X$ be an $n$-simplicial object of an additive category $\mathcal A$.
Put $\widetilde X_{\mathbf p}=X_{\mathbf p}$ for $\mathbf p\geq0$ and
zero otherwise, and let the $i$th differential be the alternating sum of
the faces in the $i$th variable. Then $\widetilde X$ is an $n$-complex,
and every total degree has only finitely many summands. If $\mathcal A$ is
abelian, this construction is an exact additive functor in each simplicial
variable and its signed total is exact.
\end{lemma}

\begin{proof}
The simplicial identities give $d_i^2=0$ in each variable. Faces in
different variables commute, so the $d_i$ commute. In total degree $r$ the
indices are $n$-tuples of nonnegative integers summing to $r$, a finite set;
thus only finite biproducts are used. Exactness follows degreewise from
exactness of finite biproducts in an abelian category and from the usual
alternating-face complex in each variable (Stacks Project, Tags 0194 and
0195). The final assertion is the resulting finite total complex.
\end{proof}

\begin{definition}[Shuffle and Alexander--Whitney maps]
\label{definition-illusie-I-shuffle-aw}
For a bisimplicial object $X$ in an additive category, write
$\operatorname{Tot}\widetilde X$ for the signed total of its two chain
directions and $\widetilde{\Delta X}$ for the chain complex of its
diagonal. For $p+q=r$, let a $(p,q)$-shuffle be a partition of
$\{1,\ldots,r\}$ into increasing lists $\mu$ and $\nu$ of lengths $p$
and $q$, and let $\epsilon(\mu,\nu)$ be its inversion parity. Define
\[
 \operatorname{Sh}_{p,q}=
 \sum_{(\mu,\nu)}(-1)^{\epsilon(\mu,\nu)}
 X(s^{\nu_q}\cdots s^{\nu_1},s^{\mu_p}\cdots s^{\mu_1}) :
 X_{p,q}\longrightarrow X_{r,r}.
\]
The shuffle map is the direct-sum map $\operatorname{Sh}$ with these
components. Define $\operatorname{AW}$ on the $r$th diagonal term by its
$X_{p,q}$ component ($p+q=r$)
\[
 X(d^{r}\cdots d^{p+1},d^0\cdots d^0):X_{r,r}\longrightarrow X_{p,q}.
\]
\end{definition}

\begin{theorem}[Eilenberg--Zilber--Cartier]
\label{theorem-illusie-I-eilenberg-zilber}
For every bisimplicial object $X$ of an additive category, the maps in
Definition \ref{definition-illusie-I-shuffle-aw} are natural chain maps,
$\operatorname{AW}\,\operatorname{Sh}=1$, and
$\operatorname{Sh}\,\operatorname{AW}$
is naturally chain-homotopic to $1$. Both maps are the identity in degree
zero. The shuffle maps are associative for trisimplicial objects and the
Alexander--Whitney maps are coassociative: the two composites obtained by
first combining the first two or the last two variables are equal.
\end{theorem}

\begin{proof}
The simplicial identities move a face past a degeneracy. In the signed sum
defining $\operatorname{Sh}$, the terms in which a face meets a degeneracy
cancel in pairs; the uncancelled terms are precisely the alternating-face
differential on the source and target. This proves that $\operatorname{Sh}$
is a chain map. The same face calculation, without degeneracies, proves
that $\operatorname{AW}$ is a chain map. In the composite
$\operatorname{AW}\operatorname{Sh}$ all non-identity terms contain a face
and degeneracy in the wrong order and cancel in the shuffle pairing; the
identity term has coefficient one.

For completeness, filter the endomorphism $\operatorname{Sh}\operatorname{AW}-1$
by the first position at which a diagonal simplex is not in front--back
normal form. The simplicial identities give a canonical degeneracy at that
position. Sending a summand to this degeneracy, with the sign prescribed by
the number of preceding faces, defines maps $H_r:X_{r,r}\to X_{r+1,r+1}$.
The same pairwise cancellation gives
\[
 dH+Hd=\operatorname{Sh}\operatorname{AW}-1.
\]
This is the standard prism contraction; it is natural because it uses only
faces, degeneracies and the ordered first exceptional position. Thus the
homotopy is functorial. For a trisimplicial object, a shuffle is a choice of
an ordered partition of one finite ordinal into three increasing lists.
Both iterated shuffles sum once over exactly these partitions with the same
inversion sign, proving associativity. The dual front--back face strings in
the two iterated Alexander--Whitney composites are identical, proving
coassociativity. In degree zero there is one empty shuffle and no face, so
both maps are the identity.
\end{proof}

\begin{theorem}[Iterated Dold--Kan normalization]
\label{theorem-illusie-I-iterated-dold-kan}
Let $\mathcal A$ be an abelian category and let $n\geq1$. Applying the
normal-complex functor in each of the $n$ simplicial variables gives a
functor from $n$-simplicial objects to $n$-fold chain complexes. In each
variable the normal and Dold--Puppe functors are quasi-inverse equivalences.
Consequently, applying them successively in any order gives canonically
isomorphic $n$-fold functors, and their signed total complexes are canonically
isomorphic after any grouping of variables. These isomorphisms are natural
and preserve the alternating-face differentials.

\end{theorem}

\begin{proof}
For one variable this is the Dold--Kan equivalence (Stacks, Tag 019G),
with the normal-complex convention translated by the canonical sign and
reversal isomorphism described in Section I.1.3. For two variables, apply
that equivalence in the first variable in the functor category of simplicial
objects in the second variable. Exactness and functoriality of the
construction are degreewise, so the same argument applies in the second
variable. The two orders commute because each face or degeneracy acts in
only one variable; the two canonical composites are therefore identical on
all multidegrees and structure maps. Induction gives the assertion for $n$
variables. The units and counits of the one-variable equivalences provide
quasi-inverse comparison maps in every variable.

The signed total of the resulting $n$-complex is the construction in
Definition \ref{definition-illusie-I-signed-total}. Grouping two adjacent
variables first changes the differential by the prefix sign for the grouped
block and then by the inner prefix sign. Their sum is the same prefix sign
as in the ungrouped total, by Lemma
\ref{lemma-illusie-I-signed-total-differential}. Thus the canonical
multivariable comparisons commute with totalization and with every
regrouping. Naturality follows from naturality of the one-variable units,
counits and face maps.
\end{proof}
